\documentclass[letterpaper]{article}
\PassOptionsToPackage{unicode=true}{hyperref} % options for packages loaded elsewhere
\PassOptionsToPackage{hyphens}{url}
\usepackage{graphicx} % Required for inserting images


\def\tightlist{}


\graphicspath{ {../images/} }
\usepackage[margin=1in]{geometry}
\usepackage{tikz}
\usepackage[dvipsnames]{xcolor}
\usetikzlibrary{calc}




\usepackage[utf8]{inputenc}
\usepackage[T1]{fontenc}
\usepackage{lmodern}
\usepackage[english]{babel}


% Useful packages
\usepackage{hyperref}
\usepackage{amsmath, amssymb}
\usepackage{fancyhdr}
\usepackage{setspace}
\usepackage{tocloft}
\usepackage{titlesec}


\usepackage{unicode-math}



% Code highlighting (Pandoc uses listings or minted, but listings is safer without shell escape)
\usepackage{listings}
\lstset{
  basicstyle=\ttfamily\small,
  breaklines=true,
  frame=single,
  backgroundcolor=\color[gray]{0.95}
}



%Modifiable Commands; defaults set to none
\newcommand{\titleDOC}{}
\newcommand{\authorDOC}{Daniel Sabalakov}
\newcommand{\dateDOC}{\today}
\newcommand{\classDOC}{}
\newcommand{\emailDOC}{danielsabalakov@gmail.com}
% DEFAULT QUOTE IS BY ISSAC NEWTON. IF YOU DON'T WANT IT, DO: quote: @none
\newcommand{\quoteDOC}{``\textit{Truth is ever to be found in the simplicity, and not in the multiplicity and confusion of things}'' - Issac Newton}

% Header/Footer
\pagestyle{fancy}
\fancyhf{}
\rhead{\thepage}
\lhead{\leftmark}





% %Title Arguments
% \title{\TITLEOFDOCUMENT}
% \author{\AUTHOROFDOCUMENT}
% \date{\today}


%
% ANYTHING BELOW HERE CAN GET PROCEDURALLY GENERATED
%
% BEGINNING OF DOCUMENT
%
%
%
\begin{document}

%titlepage








\renewcommand{\authorDOC}{Daniel Sabalakov}

\renewcommand{\classDOC}{DE Intro to Statistics I}

\renewcommand{\titleDOC}{What's the Tea}

\renewcommand{\quoteDOC}{``\textit{It's times like these you learn to live again + It's times like these you give and give again + It's times like these you learn to love again + It's times like these time and time again }'' - Foo Fighters}





\begin{titlepage}
  \titleDOC
  \classDOC
  \dateDOC
  \authorDOC
\end{titlepage}

\section{What\textquotesingle s the Tea?}\label{whats-the-tea}

Researchers investigated the possible beneficial effect on heart health of drinking black tea and whether adding milk to tea reduces any possible benefit. Twenty-four volunteers were randomly assigned to one of three groups. Every day for a month, participants in group 1 drank two cups of hot black tea without milk, participants in group 2 drank two cups of hot black tea with milk, and participants in group 3 drank two cups of hot water but no tea. At the end of the month, the researchers measured the change in each of the participants' heart health.

\begin{enumerate}
\def\labelenumi{(\alph{enumi})}
\tightlist
\item
  Did the researchers conduct an experiment or an observational study? Explain. Cite at least four significant reasons in your answer.
\end{enumerate}

\begin{itemize}
\tightlist
\item
  The researchers conducted an experiment, NOT an observational study. This is because of the following four reasons:

  \begin{itemize}
  \tightlist
  \item
    Randomization: Volunteers are "randomly" selected to one of three groups. This follows the principle of
  \item
    Replication:
  \item
    Control:
  \item
    Presence of variables:
  \end{itemize}
\end{itemize}

\begin{enumerate}
\def\labelenumi{(\alph{enumi})}
\setcounter{enumi}{1}
\item
  Why did the researchers include a group who drank hot water but no tea? Explain.
\item
  Is it reasonable to generalize the results of the study beyond the 24 participants? Explain.
\end{enumerate}

The researchers randomly assigned participants to treatment groups.

\begin{enumerate}
\def\labelenumi{(\alph{enumi})}
\setcounter{enumi}{3}
\item
  What was the reasoning behind that? Explain.
\item
  Describe a method for randomly assigning the 24 volunteers to the three treatment groups. Assume that the researchers wanted an equal number of subjects in each group.
\item
  One aspect of heart health the researchers measured was improvement in blood pressure. The results of the study are graphed below. Is there an association with one or more of the treatment groups and improvement in blood pressure? Explain in 2-3 paragraphs.
\end{enumerate}



\end{document}
